\section{Sheaf Database}
\label{sec:sheaf-database}

The Sheaf Database (SheafDB) is the primary contribution of this paper:
a storage and query system based on sheaf theory.

\subsection{Sheaf Storage}
\label{sec:sheaf:storage}

Facts are stored directly as sections over a poset of contexts, without
decomposition. Each fact is assigned to its maximal context and linked to
ancestor contexts via restriction maps. The storage layer maintains:

\begin{itemize}
    \item \textbf{Context index:} maps a context to its set of sections.
    \item \textbf{Stalk index:} maps an entity to contexts mentioning it.
    \item \textbf{Restriction index:} maps $(c_{\text{broad}}, c_{\text{narrow}})$
          to the corresponding restriction map.
    \item \textbf{Neighbourhood index:} maps a context to its parents/children.
    \item \textbf{Global section cache:} cached result of gluing local sections.
\end{itemize}

\subsection{Restriction Maps}
\label{sec:sheaf:restriction}

For $c \leq d$, the restriction map $\rho_{d,c}$ returns the fact valid
in context $d$ as seen from sub-context $c$. Since facts are context-aware,
restriction may filter object slots not meaningful in $c$.

\subsection{Global Sections}
\label{sec:sheaf:global}

Global section computation reconstructs facts valid in the root context
from compatible local sections. The algorithm processes the context poset
bottom-up, checking pairwise compatibility and gluing upward.

\subsection{Query Execution}
\label{sec:sheaf:query}

Context-scoped queries examine only the relevant context's sections,
avoiding the join overhead of the KG. Global queries require gluing
compatible local sections, which scales with the number of contexts
rather than triples.
